\chapter{Về Sai Lầm}
\markboth{Về Sai Lầm}{On Mistakes}

\section{Bài sai là bài học đang chờ em hiểu.}
\begin{truyen}{Bài sai đang chờ em}{The Wrong Answer Is Waiting for You}
\chuhoa{T}{hầy nói khi trả bài: ``Đừng vội gấp bài sai bỏ đi. Mỗi bài sai là một bài học đang chờ em — em đọc lại đề, xem em sai ở bước nào, vì sao sai, rồi sửa. Khi em hiểu rồi, bài sai ấy mới trở thành bài học thật. Còn nếu em bỏ qua, nó vẫn sẽ quay lại trong bài sau.''}
\textit{(When returning papers the teacher said: ``Don't rush to fold the wrong ones away. Each wrong answer is a lesson waiting for you — read the question again, see where you went wrong and why, then fix it. When you understand, that wrong answer becomes a real lesson. If you skip it, it will come back in the next test.'')}
\end{truyen}

\begin{baihoc}
Bài sai chỉ trở thành bài học khi em dừng lại, xem lại và hiểu vì sao sai.
\textit{(A wrong answer only becomes a lesson when you stop, look back, and understand why it was wrong.)}
\end{baihoc}

\begin{ghinhoanh}
\textbf{Short English to remember:}

A wrong answer is a lesson waiting for you to understand it.

\end{ghinhoanh}

\begin{tuhocnhanh}
1. Em có thường xem lại bài sai và tìm hiểu vì sao không? \textit{Do you usually review wrong answers and find out why?}

2. Bài sai nào gần đây em đã ``hiểu'' và sửa? \textit{Which recent wrong answer have you understood and fixed?}
\end{tuhocnhanh}
\ngancach

\section{Người không bao giờ sai thường là người không làm gì mới.}
\begin{truyen}{Không sai thì không mới}{No Mistakes Often Means Nothing New}
\chuhoa{T}{hầy nói với em luôn làm đúng vì chỉ chọn bài dễ: ``Nếu em chỉ làm bài em chắc chắn đúng, em đang đứng yên. Người không bao giờ sai thường là người không thử cái mới — không làm bài khó, không phát biểu khi chưa chắc. Sai là dấu hiệu em đang bước ra vùng mới.''}
\textit{(The teacher told a student who always got things right by only doing easy problems: ``If you only do what you're sure to get right, you're standing still. People who never make mistakes often never try anything new — no hard problems, no speaking up when unsure. Making mistakes is a sign you're stepping into new territory.'')}
\end{truyen}

\begin{baihoc}
Sai trong khi thử cái mới là bình thường; tránh sai bằng cách không thử mới là đáng lo.
\textit{(Mistakes when trying something new are normal; avoiding mistakes by not trying is what's worrying.)}
\end{baihoc}

\begin{ghinhoanh}
\textbf{Short English to remember:}

Those who never make mistakes are often those who never try anything new.

\end{ghinhoanh}

\begin{tuhocnhanh}
1. Em có đang tránh sai bằng cách chỉ làm bài dễ không? \textit{Are you avoiding mistakes by only doing easy work?}

2. Thử ``cái mới'' trong học với em là gì? \textit{What is ``something new'' for you in learning?}
\end{tuhocnhanh}
\ngancach

\section{Sai một lần mà hiểu được, còn quý hơn đúng mà không hiểu.}
\begin{truyen}{Sai mà hiểu còn quý hơn đúng mà không hiểu}{Wrong But Understanding Is Worth More}
\chuhoa{M}{ột em làm sai nhưng sau đó hỏi thầy và ghi lại cách đúng. Thầy nói trước lớp: ``Em này sai nhưng đã hiểu vì sao và cách sửa. Có em làm đúng nhưng nhìn lại không giải thích được — đúng kiểu may. Sai một lần mà hiểu được còn quý hơn đúng mà không hiểu. Vì lần sau em ấy sẽ không sai nữa.''}
\textit{(A student got it wrong but then asked the teacher and wrote down the correct method. The teacher told the class: ``This one was wrong but understood why and how to fix it. Some get it right but can't explain — luck. Wrong once but understanding is worth more than right without understanding. Next time they won't make that mistake again.'')}
\end{truyen}

\begin{baihoc}
Đúng mà không hiểu chỉ là may một lần; sai mà hiểu là nền cho đúng mãi sau.
\textit{(Right without understanding is luck once; wrong but understanding is the base for being right from then on.)}
\end{baihoc}

\begin{ghinhoanh}
\textbf{Short English to remember:}

Being wrong once but understanding is worth more than being right without understanding.

\end{ghinhoanh}

\begin{tuhocnhanh}
1. Em có từng làm đúng nhưng không giải thích được không? \textit{Have you ever been right but couldn't explain?}

2. Sau khi sai, em có tìm hiểu cho đến khi hiểu không? \textit{After a mistake, do you dig until you understand?}
\end{tuhocnhanh}
\ngancach

\section{Sai lầm là giá học phí của kinh nghiệm.}
\begin{truyen}{Giá học phí của kinh nghiệm}{The Tuition of Experience}
\chuhoa{T}{hầy nói: ``Mỗi lần em sai — sai bài, sai lời nói, sai cách cư xử — em đang trả một khoản học phí. Khoản đó không lãng phí nếu em nhận ra mình sai ở đâu và không lặp lại. Sai lầm là giá học phí của kinh nghiệm. Đắt hay rẻ tùy em có học từ nó hay không.''}
\textit{(The teacher said: ``Every time you're wrong — wrong answer, wrong word, wrong behavior — you're paying tuition. That payment isn't wasted if you see where you went wrong and don't repeat it. Mistakes are the tuition of experience. Whether that's expensive or cheap depends on whether you learn from them.'')}
\end{truyen}

\begin{baihoc}
Sai lầm chỉ ''đắt'' khi em trả xong mà vẫn không học; học được thì nó thành đầu tư.
\textit{(Mistakes are only ``expensive'' when you pay and still don't learn; if you learn, they become an investment.)}
\end{baihoc}

\begin{ghinhoanh}
\textbf{Short English to remember:}

Mistakes are the tuition fee of experience.

\end{ghinhoanh}

\begin{tuhocnhanh}
1. Em có xem sai lầm là ''học phí'' hay chỉ là thất bại? \textit{Do you see mistakes as ``tuition'' or just failure?}

2. Sai lầm gần đây em đã ''trả học phí'' và học được gì? \textit{What recent mistake did you ``pay tuition'' for and what did you learn?}
\end{tuhocnhanh}
\ngancach

\section{Người thông minh học từ sai lầm của mình, người khôn ngoan học từ sai lầm của người khác.}
\begin{truyen}{Học từ sai lầm của người khác}{Learning from Others' Mistakes}
\chuhoa{T}{hầy kể: ``Khi bạn em bị trừ điểm vì quên làm bài, em có thể nghĩ: mình sẽ không quên. Khi bạn bị la vì nói chuyện trong giờ, em có thể nhớ: mình sẽ giữ trật tự. Đó là học từ sai lầm của người khác — không cần trả giá bằng chính mình. Người thông minh học từ sai của mình; người khôn ngoan còn học từ sai của người khác.''}
\textit{(The teacher said: ``When your friend loses points for forgetting homework, you can think: I won't forget. When a friend is scolded for talking in class, you can remember: I'll stay quiet. That's learning from others' mistakes — without paying the price yourself. The smart learn from their own; the wise also learn from others'.'')}
\end{truyen}

\begin{baihoc}
Học từ sai của mình là cần; học thêm từ sai của người khác là khôn — đỡ trả giá.
\textit{(Learning from your own mistakes is necessary; learning from others' is wise — you pay less.)}
\end{baihoc}

\begin{ghinhoanh}
\textbf{Short English to remember:}

The smart learn from their mistakes; the wise learn from others' mistakes too.

\end{ghinhoanh}

\begin{tuhocnhanh}
1. Em có quan sát và rút ra bài học từ sai lầm của bạn không? \textit{Do you observe and learn from friends' mistakes?}

2. Sai lầm nào của người khác em đã ''học giùm''? \textit{Which mistake of someone else have you ``learned for free''?}
\end{tuhocnhanh}
\ngancach

\section{Đừng xấu hổ vì sai, hãy xấu hổ nếu không chịu sửa.}
\begin{truyen}{Xấu hổ khi không sửa}{Ashamed If You Don't Fix It}
\chuhoa{T}{hầy nói với em hay giấu bài sai: ``Sai thì ai cũng sai. Không có gì phải xấu hổ khi làm sai bài. Nhưng nếu em biết sai rồi mà không chịu xem lại, không hỏi, không sửa — đó mới đáng xấu hổ. Vì em đang bỏ qua cơ hội tiến lên. Đừng xấu hổ vì sai; hãy xấu hổ nếu không chịu sửa.''}
\textit{(The teacher told a student who often hid wrong papers: ``Everyone makes mistakes. There's nothing to be ashamed of in a wrong answer. But if you know it's wrong and won't look again, ask, or fix it — that's what's shameful. You're throwing away the chance to improve. Don't be ashamed of being wrong; be ashamed of not fixing it.'')}
\end{truyen}

\begin{baihoc}
Sai không đáng xấu hổ; đáng xấu hổ là biết sai mà không sửa.
\textit{(Being wrong isn't shameful; knowing you're wrong and not fixing it is.)}
\end{baihoc}

\begin{ghinhoanh}
\textbf{Short English to remember:}

Don't be ashamed of being wrong; be ashamed if you don't fix it.

\end{ghinhoanh}

\begin{tuhocnhanh}
1. Em có hay giấu bài sai hoặc không muốn xem lại không? \textit{Do you often hide wrong work or avoid looking at it?}

2. Sau khi biết mình sai, em có chủ động sửa không? \textit{After knowing you're wrong, do you actively fix it?}
\end{tuhocnhanh}
\ngancach

\section{Mỗi lỗi sai đều đang chỉ đường cho em đi đúng hơn.}
\begin{truyen}{Lỗi sai chỉ đường}{Mistakes Point the Way}
\chuhoa{T}{hầy giải thích: ``Mỗi lỗi sai giống như biển báo — nó báo em đã đi lệch chỗ này. Nếu em đọc được biển báo ấy, em sẽ biết chỗ nào cần chỉnh. Mỗi lỗi sai đều đang chỉ đường cho em đi đúng hơn. Vấn đề không phải là không sai, mà là sai xong em có đọc biển báo không.''}
\textit{(The teacher explained: ``Each mistake is like a sign — it says you went wrong here. If you read that sign, you know what to correct. Every mistake is pointing you toward a better path. The issue isn't never making mistakes; it's whether you read the sign after you're wrong.'')}
\end{truyen}

\begin{baihoc}
Sai lầm là tín hiệu chỉ chỗ cần sửa; bỏ qua tín hiệu thì sẽ sai lại.
\textit{(Mistakes are signals showing what to fix; ignoring the signal means erring again.)}
\end{baihoc}

\begin{ghinhoanh}
\textbf{Short English to remember:}

Every mistake is pointing you toward the right way.

\end{ghinhoanh}

\begin{tuhocnhanh}
1. Em có xem lỗi sai như ''biển báo'' hay như thất bại? \textit{Do you see mistakes as ``signposts'' or as failure?}

2. Lỗi sai gần đây đã ''chỉ đường'' cho em điều gì? \textit{What did a recent mistake ``point you toward''?}
\end{tuhocnhanh}
\ngancach

\section{Người trưởng thành là người nhìn lại lỗi sai mà không đổ lỗi.}
\begin{truyen}{Nhìn lại lỗi không đổ lỗi}{Looking Back Without Blaming Others}
\chuhoa{T}{hầy nói khi hai em cãi nhau vì nhóm làm bài sai: ``Người trưởng thành khi nhìn lại lỗi sai sẽ hỏi: mình đã làm gì sai, mình có thể sửa thế nào. Người chưa trưởng thành sẽ đổ cho bạn, cho thầy, cho đề khó. Đổ lỗi không sửa được gì; nhìn vào mình mới sửa được. Hãy tập nhìn lại lỗi mà không đổ lỗi.''}
\textit{(When two students argued over a wrong group assignment the teacher said: ``When looking back at a mistake, a mature person asks: what did I do wrong, how can I fix it. An immature one blames the friend, the teacher, the hard question. Blaming fixes nothing; looking at yourself does. Practice looking back at mistakes without blaming others.'')}
\end{truyen}

\begin{baihoc}
Đổ lỗi giữ em đứng yên; nhận phần mình và sửa mới đưa em tiến.
\textit{(Blaming keeps you stuck; owning your part and fixing it moves you forward.)}
\end{baihoc}

\begin{ghinhoanh}
\textbf{Short English to remember:}

Growing up means looking at your mistakes without blaming others.

\end{ghinhoanh}

\begin{tuhocnhanh}
1. Khi sai, em thường đổ lỗi hay nhìn vào phần mình? \textit{When you're wrong, do you usually blame or look at your part?}

2. Đổ lỗi có giúp em sửa được không? \textit{Does blaming help you fix anything?}
\end{tuhocnhanh}
\ngancach

\section{Sai không đáng sợ, lặp lại sai mới đáng sợ.}
\begin{truyen}{Lặp lại sai mới đáng sợ}{Repeating the Same Mistake Is What's Scary}
\chuhoa{T}{hầy nói với em lần thứ ba quên cùng một dạng bài: ``Sai lần đầu là bình thường. Sai lần hai có thể do chưa kịp ôn. Nhưng sai cùng một chỗ lần thứ ba nghĩa là em chưa thật sự xem lại và sửa. Sai không đáng sợ; lặp lại sai mới đáng sợ — vì đó là em không học.''}
\textit{(The teacher told a student who had forgotten the same type of problem a third time: ``Wrong the first time is normal. Wrong the second might be not enough review. But wrong the same way a third time means you haven't really looked back and fixed it. Being wrong isn't scary; repeating the same mistake is — because that means you're not learning.'')}
\end{truyen}

\begin{baihoc}
Sai một lần là bài học; sai y hệt lần nữa là bỏ qua bài học.
\textit{(Wrong once is a lesson; wrong the same way again is skipping the lesson.)}
\end{baihoc}

\begin{ghinhoanh}
\textbf{Short English to remember:}

Being wrong isn't scary; repeating the same mistake is.

\end{ghinhoanh}

\begin{tuhocnhanh}
1. Em có từng sai cùng một kiểu nhiều lần không? \textit{Have you ever made the same kind of mistake many times?}

2. Làm sao để không lặp lại sai? \textit{How can you avoid repeating the same mistake?}
\end{tuhocnhanh}
\ngancach

\section{Nếu em chưa từng sai, có thể em chưa thử đủ nhiều.}
\begin{truyen}{Chưa sai có thể là chưa thử đủ}{Never Wrong May Mean Not Trying Enough}
\chuhoa{T}{hầy nói với em luôn an toàn, chỉ làm bài trong khả năng: ``Nếu em chưa từng sai, có thể em đang chơi quá an toàn — chỉ làm bài em chắc đúng, không thử bài khó, không phát biểu khi chưa chắc. Người tiến nhanh là người thử nhiều, sai nhiều, và sửa nhiều. Em hãy thử mở rộng vùng an toàn một chút.''}
\textit{(The teacher told a student who always played it safe: ``If you've never been wrong, you might be playing too safe — only doing what you're sure to get right, not trying hard problems, not speaking when unsure. Those who improve fast try a lot, err a lot, and fix a lot. Try widening your comfort zone a little.'')}
\end{truyen}

\begin{baihoc}
Không sai có thể là đang tránh thử; thử nhiều thì sẽ có sai — và có tiến.
\textit{(No mistakes might mean you're avoiding trying; trying more means some wrongs — and progress.)}
\end{baihoc}

\begin{ghinhoanh}
\textbf{Short English to remember:}

If you've never been wrong, you may not have tried enough.

\end{ghinhoanh}

\begin{tuhocnhanh}
1. Em có đang ''an toàn'' quá — chỉ làm bài dễ, không thử bài khó? \textit{Are you playing it too safe — only easy work, no hard tries?}

2. Em có dám thử điều khó hơn một chút không? \textit{Do you dare to try something a bit harder?}
\end{tuhocnhanh}
\ngancach
